\subsection{Causality in \GRSync}
\label{sec:semantics:grsync:causality}

This section gives formal definition to the causality of \GRSync constructs.
%
In everyday language, causality refers to the cause-and-effect relationship between several
phenomena: it is the influence by which one phenomenon (the cause) contributes to the production of
another phenomenon (the effect), with the same cause producing the same effects under the same
conditions.
%
Similarly, in synchronous programming, we say that a component is \emph{causal} if the same input
flows produce the same output flows.

Causality verification is performed statically at compile time. It involves analyzing the
relationships between defined identifiers (formalized in \Cref{sec:semantics:grsync:causality:defs}) and their
dependencies to construct a dependency order (defined in \Cref{sec:semantics:grsync:causality:dependencies}).
A component, and more generally a block of equations, is causal if no identifier depends on itself
(as formalized in \Cref{sec:semantics:grsync:causality:causal}).
%
From this dependency order, we can classify identifiers based on the maximum length of their
dependency chains. This classification is referred to as the \emph{rank}
(defined in \Cref{sec:semantics:grsync:causality:rank}). Using ranks, we enable an inductive
reasoning on identifiers defined by a causal block of equations (\Cref{sec:semantics:grsync:causality:induction})
and derive a scheduling (\Cref{sec:semantics:grsync:causality:scheduling}) in which equations can
be computed without any deadlock caused by circular dependencies.

\subsubsection{Defined identifiers}
\label{sec:semantics:grsync:causality:defs}

\newcommand{\defs}[1]{{\ensuremath{\textit{defs}\left(#1\right)}}}

Equations, patterns, and event patterns introduce identifiers into scope. To reason formally about
causality, we define the function $\defs{\epsilon}$, which returns the set of identifiers defined
by a given \GRSync construct $\epsilon$ (an equation, a pattern, or an event pattern).

\begin{definition}[Definition set]
    Let $\epsilon$ be an element of \GRSync. The set of identifiers defined by $\epsilon$ is
    denoted \defs{\epsilon}, and its definition is given in
    \Crefrange{rule:semantics:grsync:causality:defs:eq}%
    {rule:semantics:grsync:causality:defs:epat:edge}.
\end{definition}

\paragraph{Equations}

A block of equations defines the union of those defined by each equation (\Cref{rule:semantics:%
    grsync:causality:defs:eq}). A declaration $x^+ = e$ defines the identifiers on its left-hand
side, which are $x^+$ (\Cref{rule:semantics:grsync:causality:defs:eq:decl}). In \kwf{match}
equations, the defined identifiers are the $\kwf{vars}\{x^+\}$ (\Cref{rule:semantics:grsync:%
    causality:defs:eq:match}). In \kwf{when} equations, the defined identifiers are the
$\kwf{vars}\{x^+\}$ plus the memories of the signals defined in \init{x_m^+ = c^+} (\Cref{rule:%
    semantics:grsync:causality:defs:eq:when}).

\begin{align}
    \defs{eq^+}                                                           & = \bigsqcup \defs{eq}^+ \label{rule:semantics:grsync:causality:defs:eq}              \\
    \defs{x^+ = e}                                                        & = \{x^+\} \label{rule:semantics:grsync:causality:defs:eq:decl}                       \\
    \defs{\matchTab{e}{x^+}{(\branch{pat_i^+}{e_i}{eq_i^+})^+}}           & = \{x^+\} \label{rule:semantics:grsync:causality:defs:eq:match}                      \\
    \defs{\whenTab{x^+}{x_m^+ = c^+}{(\branch{epat_i^+}{e_i}{eq_i^+})^+}} & = \{x^+\} \cup \{\last{x_m}^+\} \label{rule:semantics:grsync:causality:defs:eq:when}
\end{align}

\paragraph{Patterns}

Patterns define identifiers when matching values. A tuple pattern $pat^+$ defines the union of its
sub-patterns (\Cref{rule:semantics:grsync:causality:defs:pat}). A constant $c$ or a wildcard
\verb|_| defines no identifier (\Crefrange{rule:semantics:grsync:causality:defs:pat:const}%
{rule:semantics:grsync:causality:defs:pat:wild}). An identifier pattern $x$ defines the
identifier $x$ (\Cref{rule:semantics:grsync:causality:defs:pat:ident}).

\begin{align}
    \defs{pat^+} & = \bigsqcup \defs{pat}^+ \label{rule:semantics:grsync:causality:defs:pat} \\
    \defs{x}     & = \{x\} \label{rule:semantics:grsync:causality:defs:pat:ident}            \\
    \defs{c}     & = \emptyset \label{rule:semantics:grsync:causality:defs:pat:const}        \\
    \defs{\_}    & = \emptyset \label{rule:semantics:grsync:causality:defs:pat:wild}
\end{align}

\paragraph{Event patterns}

Event patterns define identifiers upon event detection. A tuple $epat^+$ defines the union of its
components (\Cref{rule:semantics:grsync:causality:defs:epat}). An event detection \detect{x}{y} defines
$x$ (\Cref{rule:semantics:grsync:causality:defs:epat:event}), while a rising edge $e$ defines
no identifier (\Cref{rule:semantics:grsync:causality:defs:epat:edge}).

\begin{align}
    \defs{epat^+}        & = \bigsqcup \defs{epat}^+ \label{rule:semantics:grsync:causality:defs:epat} \\
    \defs{\detect{x}{y}} & = \{x\} \label{rule:semantics:grsync:causality:defs:epat:event}             \\
    \defs{e}             & = \emptyset \label{rule:semantics:grsync:causality:defs:epat:edge}
\end{align}

\subsubsection{Dependencies}
\label{sec:semantics:grsync:causality:dependencies}

Dependencies express how a computation---or the definition of an identifier---may rely on the value
of other identifiers. They form the foundation for scheduling and for detecting causality
violations. This section defines how to extract dependencies from \GRSync constructs and how to
derive the dependency relation between identifiers.

\paragraph{Dependency set of constructs}
\label{sec:semantics:grsync:causality:dependencies:deps}

\newcommand{\deps}[1]{{\ensuremath{\textit{deps}\left(#1\right)}}}

Before defining a dependency order between identifiers, we first need to define what are the
identifiers that need to be available before performing a computation. Those identifiers are
collected in a dependency set.

\begin{definition}[Dependency set of constructs]
    Let $\epsilon$ be an element of \GRSync (an expression, an equation, or an event pattern).
    The set of dependencies of $\epsilon$ is denoted \deps{\epsilon}, and its definition is given in
    \Crefrange{rule:semantics:grsync:causality:dependencies:deps:e:const}%
    {rule:semantics:grsync:causality:dependencies:deps:epat:edge}.
\end{definition}

\subparagraph{Expressions}

The dependencies of an expression are the identifiers that are needed for its evaluation.
%
A constant expression $c$ introduces no dependencies, as its value is always accessible.
%
An identifier usage $x$ introduces a dependency on the identifier $x$, reflecting that the current
value of $x$ is required.
%
A memory access expression of the form $\last{x}$ depends on the previous value of $x$,
represented as a dependency on $\last{x}$.
%
An event emission expression \emit{e} depends on the same identifiers as the expression $e$
itself, since $e$ must be evaluated to determine the emitted value.
%
An operation $op(e^+)$ depends on the union of the dependencies of the sub-expressions
$e^+$, each argument $e$ must be computable before the operator $op$ can be applied.
%
Finally, a component application $f(e^+)$ similarly depends on the union of the dependencies
of its arguments $e^+$, since the arguments must be evaluated prior to the call.

\begin{align}
    \deps{c}        & = \emptyset \label{rule:semantics:grsync:causality:dependencies:deps:e:const}        \\
    \deps{x}        & = \{x\} \label{rule:semantics:grsync:causality:dependencies:deps:e:ident}            \\
    \deps{\last{x}} & = \{\last{x}\} \label{rule:semantics:grsync:causality:dependencies:deps:e:last}      \\
    \deps{\emit{e}} & = \deps{e} \label{rule:semantics:grsync:causality:dependencies:deps:e:emit}          \\
    \deps{op(e^+)}  & = \bigcup \deps{e}^+ \label{rule:semantics:grsync:causality:dependencies:deps:e:op}  \\
    \deps{f(e^+)}   & = \bigcup \deps{e}^+ \label{rule:semantics:grsync:causality:dependencies:deps:e:app}
\end{align}

\subparagraph{Equations}

The dependencies of equations are the identifiers required to be available in order to compute the
equation.
%
A block of equations $eq^+$ depends on the union of the dependencies of each individual equation in
the block, minus the identifiers they define, as they are computed by the equations and then not
pre-required.
%
In a declaration of the form $x^+ = e$, the dependencies are exactly those of the expression $e$, as
$e$ must be evaluated to define the values of $x^+$.
%
The dependencies of a \kwf{match} equation of the form \match{e}{x^+}{
    (\branch{pat_i^+}{e_i}{eq_i^+})^+} come from different sources. First, it include the
dependencies of the matched expression $e$. Then, each branch introduces two additional dependency
sources: the guard $e_i$ and the corresponding equation block $eq_i^+$.
%
A \kwf{when} equation of the form \when{x^+}{x_m^+ = c^+}{(\branch{epat_i^+}{e_i}{eq_i^+})^+}
similarly depends on the contents of each of its branches. For each branch, the event pattern
$epat_i^+$ may introduce dependencies (\eg for event detection), the guard expression $e_i$
introduces other dependencies, and the associated block of equations $eq_i^+$ also contributes.
But the block of equations $eq_i^+$ represent flow updates, so all the identifiers defined by the
\kwf{when} equation have to be removed, along the memories of defined signals ($\last{x_m}^+$). It
was not necessary in \kwf{match} equations since $x^+$ are all defined in each branch.

\begin{align}
    \deps{eq^+}                                                                      & =
    \left(\bigcup \deps{eq}^+\right) \backslash \defs{eq^+}
    \label{rule:semantics:grsync:causality:dependencies:deps:eq}                         \\
    \deps{x^+ = e}                                                                   & =
    \deps{e}
    \label{rule:semantics:grsync:causality:dependencies:deps:eq:decl}                    \\
    \deps{\mathtiny{\matchTab{e}{x^+}{(\branch{pat_i^+}{e_i}{eq_i^+})^+}}}           & =
    \deps{e} \cup \bigcup \left(\deps{e_i} \cup \deps{eq_i^+} \backslash \defs{pat_i^+}\right)^+
    \label{rule:semantics:grsync:causality:dependencies:deps:eq:match}                   \\
    \deps{\mathtiny{\whenTab{x^+}{x_m^+ = c^+}{(\branch{epat_i^+}{e_i}{eq_i^+})^+}}} & =
    \begin{array}{l}
        \bigcup \left(\deps{epat_i^+} \cup \deps{e_i} \cup \deps{eq_i^+} \backslash \defs{epat_i^+}\right)^+ \\
        \qquad \backslash (\{x^+\} \cup \{\last{x_m}^+\})
    \end{array}
    \label{rule:semantics:grsync:causality:dependencies:deps:eq:when}
\end{align}

\subparagraph{Event patterns}

Event patterns introduce dependencies because their evaluation relies on identifier values.
%
A tuple event pattern $epat^+$ depends on the union of the dependencies of its sub-patterns.
%
A detection pattern of the form \detect{x}{y} introduces a dependency on the identifier $y$, as the event
is triggered by the value of $y$.
%
Rising edge patterns, expressions denoted by $e$, inherit their dependencies directly from the
expression $e$ itself.

\begin{align}
    \deps{epat^+}        & = \bigcup \deps{epat}^+ \label{rule:semantics:grsync:causality:dependencies:deps:epat} \\
    \deps{\detect{x}{y}} & = \{y\} \label{rule:semantics:grsync:causality:dependencies:deps:epat:event}           \\
    \deps{e}             & = \deps{e} \label{rule:semantics:grsync:causality:dependencies:deps:epat:edge}
\end{align}

\paragraph{Renaming function}
\label{sec:semantics:grsync:causality:dependencies:renaming}

Now that we have defined what are the required identifiers to be available for all the \GRSync
constructs that involve computations, we want to define a dependency relation that represents the
evaluation order for identifiers.
%
This operation is simple for declaration statements ($x^+ = e$), the declared identifiers ($x^+$)
clearly depends on the identifiers needed to compute the expression ($e$) that defines their
behavior. However, it is more complicated for \kwf{match} or \kwf{when} equations, where the
different branches implies different dependency order. In order to avoid the detection of false
causality breaches, we use a renaming technique~\cite{bourke:hal-03936656}: each branch of a
\kwf{match} or \kwf{when} equation is associated with a renaming function, which assigns unique
names to the identifiers defined within that branch, ensuring they do not conflict with other
branches.

The dependency order is usually represented as a graph where nodes are its identifiers and edges are
the dependency relation between two nodes. The dependency graphs of causal components are acyclic:
a cycle implies that an identifier depends on itself.

For example, consider the \GRSync \kwf{match} equation shown in
\Cref{fig:semantics:grsync:causality:dependencies:renaming:example:code}.
In this equation, the identifier \idf{x} depends on \idf{y} when \idf{ck} is \true, and
conversely, \idf{y} depends on \idf{x} when \idf{ck} is \false. This is perfectly valid, as
\idf{x} and \idf{y} can be computed in both branches.
%
However, if no renaming is applied during causality analysis, the global dependency order will
incorrectly merge the branch-local orders. This produces a circular dependency:
`\idf{y} depends on \idf{x}, which depends on \idf{y}`, leading to a false causality breach.

\begin{figure}[!ht]
    \centering
    \begin{subfigure}[b]{0.3\textwidth}
        \begin{lstlisting}[language=GRust]
match ck {
    true  => {
        let x: int = y;
        let y: int = 0;
    }
    false => {
        let x: int = 0;
        let y: int = x;
    }
}
\end{lstlisting}
        \caption[\GRSync \kwf{match} equation.]%
        {\GRSync \kwf{match} equation.}
        \label{fig:semantics:grsync:causality:dependencies:renaming:example:code}
    \end{subfigure}
    \quad
    \begin{subfigure}[b]{0.3\textwidth}
        \begin{align}
            \sigma_\true(\idf{x})  & = \idf{x}_\true \label{eq:semantics:grsync:causality:dependencies:renaming:example:rename:true:x}   \\
            \sigma_\true(\idf{y})  & = \idf{y}_\true \label{eq:semantics:grsync:causality:dependencies:renaming:example:rename:true:y}   \\
            \sigma_\false(\idf{x}) & = \idf{x}_\false \label{eq:semantics:grsync:causality:dependencies:renaming:example:rename:false:x} \\
            \sigma_\false(\idf{y}) & = \idf{y}_\false \label{eq:semantics:grsync:causality:dependencies:renaming:example:rename:false:y}
        \end{align}
        \caption[Renaming functions for the \kwf{match} equation.]%
        {The renaming functions.}
        \label{fig:semantics:grsync:causality:dependencies:renaming:example:rename}
    \end{subfigure}
    \quad
    \begin{subfigure}[b]{0.3\textwidth}
        \centering
        \scalebox{0.8}{\footnotesize
            \begin{tikzpicture}[node distance = 0.6cm, line width=0.5pt]
                % Place nodes
                \node (center) at (0, 0) {};
                \node[gray, ellipse, draw, dashed, above =of center,
                    minimum width=4cm, minimum height=2cm] (ellipse1) {};
                \node[gray, ellipse, draw, dashed, below =of center,
                    minimum width=4cm, minimum height=2cm] (ellipse2) {};
                \node[gray, above =of center] (ck1) {\true};
                \node[gray, below =of center] (ck2) {\false};
                \node [circle, draw, minimum width=1.2cm] (x) at (360/6 * 3:2cm) {$\idf{x}$};
                \node [circle, draw, minimum width=1.2cm] (y) at (360/6 * 6:2cm) {$\idf{y}$};
                \node [circle, draw, minimum width=1.2cm] (x1) at (360/6 * 4:2cm)
                {$\idf{x}_\true$};
                \node [circle, draw, minimum width=1.2cm] (y1) at (360/6 * 5:2cm)
                {$\idf{y}_\true$};
                \node [circle, draw, minimum width=1.2cm] (x2) at (360/6 * 2:2cm)
                {$\idf{x}_\false$};
                \node [circle, draw, minimum width=1.2cm] (y2) at (360/6 * 1:2cm)
                {$\idf{y}_\false$};

                \draw[->,>=latex] (x) to[bend right=15] (x1);
                \draw[->,>=latex] (x) to[bend left=15] (x2);
                \draw[->,>=latex] (y) to[bend left=15] (y1);
                \draw[->,>=latex] (y) to[bend right=15] (y2);
                \draw[->,>=latex] (x1) to[bend right=15] (y1);
                \draw[->,>=latex] (y2) to[bend right=15] (x2);
            \end{tikzpicture}}
        \caption[Dependency graph for the \kwf{match} equation.]%
        {Graph representing the dependency order, with branch-specific renaming.}
        \label{fig:semantics:grsync:causality:dependencies:renaming:example:graph}
    \end{subfigure}

    \caption[Renaming example with a \kwf{match} equation.]%
    {A \kwf{match} equation where each branch introduces different dependencies:
        \idf{x} depends on \idf{y} when \idf{ck} is \true, and \idf{y} depends on \idf{x}
        when \idf{ck} is \false.}
    \label{fig:semantics:grsync:causality:dependencies:renaming:example}
\end{figure}

To avoid this issue, each branch assigns fresh, unique names to the identifiers it defines. In the
compiler, these names are generated automatically. In this example, \idf{x} and \idf{y} are renamed
to $\idf{x}_\true$ and $\idf{y}_\true$ in the \true branch, and to
$\idf{x}_\false$ and $\idf{y}_\false$ in the \false branch.
%
The corresponding renaming functions, $\sigma_\true$ for the \true branch and
$\sigma_\false$ for the \false branch, are defined in
\Crefrange{eq:semantics:grsync:causality:dependencies:renaming:example:rename:true:x}%
{eq:semantics:grsync:causality:dependencies:renaming:example:rename:false:y}.
%
This transformation results in the dependency order represented by the graph shown in
\Cref{fig:semantics:grsync:causality:dependencies:renaming:example:graph}, where
$\idf{x}_\true$ depends on $\idf{y}_\true$,
$\idf{y}_\false$ depends on $\idf{x}_\false$.
%
The global identifiers \idf{x} and \idf{y} then depend on their renamed counterparts:
\idf{x} depends on $\idf{x}_\true$ and $\idf{x}_\false$,
and likewise for \idf{y}.
%
This graph cleanly separates the dependencies of each branch while preserving global relationships,
eliminating false causality violations.

\begin{definition}[Renaming function]
    \label{def:semantics:grsync:causality:dependencies:renaming}
    Let $\mathcal{V}$ be a set of identifiers.
    The function $\sigma: \mathcal{V} \rightarrow \idSet$ is a \emph{renaming function on
        $\mathcal{V}$} if for all identifiers $x, y \in \mathcal{V}$:
    $$\sigma(x) = \sigma(y) \Leftrightarrow x = y.$$
    %
    The set of all renaming functions on $\mathcal{V}$ is denoted by $\Sigma_\mathcal{V}$.

    Let $\sigma_1, \sigma_2$ be two renaming functions on a domain $\mathcal{V}$.
    $\sigma_1$ and $\sigma_2$ are \emph{disjoint on $\mathcal{S} \subseteq \mathcal{V}$}, denoted by
    $\sigma_1 \sqcap \sigma_2$, if for all identifier $x \in \mathcal{S}$:
    $$\sigma_1(x) = \sigma_2(x) \Leftrightarrow \sigma_1 = \sigma_2.$$
    %
    A set $\{\sigma^+\}$ of renaming functions on a domain $\mathcal{V}$ is disjoint on
    $\mathcal{S} \subseteq \mathcal{V}$, denoted by $\bigsqcap \sigma^+$, if all its renaming
    functions are disjoint with each other on $\mathcal{S}$.

    Let $pat^+$ be a matching pattern, $\mathcal{V}$ and $\mathcal{S}$ be sets of identifiers such that
    $\mathcal{S} \subseteq \mathcal{V}$ and, for all $x \in \mathcal{S}$, $x_{pat^+}$ is not in
    $\mathcal{V}$. We denote by $\sigma^{\mathcal{S}}_{pat^+}$ the function that maps every
    identifier $x \in \mathcal{S}$ to a labelled version $x_{pat^+}$, any other identifier
    $x \in \mathcal{V} \backslash \mathcal{S}$ is mapped to itself.
    $$\sigma^{\mathcal{S}}_{pat^+} : x \in \mathcal{V} \mapsto
        \left\{
        \begin{array}{ll}
            x_{pat^+} & \text{if} \, x \in \mathcal{S} \\
            x         & \text{otherwise}
        \end{array}
        \right.$$
\end{definition}

\begin{lemma}[Renaming function from patterns]
    \label{lem:semantics:grsync:causality:dependencies:renaming}
    Let $\mathcal{P}$ be a set of matching patterns, $\mathcal{V}$ and $\mathcal{S}$ be sets of
    identifiers such that $\mathcal{S} \subseteq \mathcal{V}$ and, for all $x \in \mathcal{S}$ and
    $pat^+ \in \mathcal{P}$, $x_{pat^+}$ is not in $\mathcal{V}$.
    %
    For all $pat^+ \in \mathcal{P}$, the function $\sigma^{\mathcal{S}}_{pat^+}$, previously defined
    in \Cref{def:semantics:grsync:causality:dependencies:renaming}, is a renaming function on $\mathcal{V}$.
    %
    The set $\{\sigma^{\mathcal{S}}_{pat^+} \mid pat^+ \in \mathcal{P}\}$, denoted by
    $\Sigma^{\mathcal{S}}_{\mathcal{P}}$, is disjoint on $\mathcal{S}$.
\end{lemma}
\begin{skippableproof}
    Let $\mathcal{P}$ be a set of matching patterns, $\mathcal{V}$ and $\mathcal{S}$ be sets of
    identifiers such that $\mathcal{S} \subseteq \mathcal{V}$ and, for all $x \in \mathcal{S}$ and
    $pat^+ \in \mathcal{P}$, $x_{pat^+}$ is not in $\mathcal{V}$.

    Let $pat^+ \in \mathcal{P}$, and $x, y \in \mathcal{V}$ such that
    $\sigma^{\mathcal{S}}_{pat^+}(x) = \sigma^{\mathcal{S}}_{pat^+}(y)$.
    For all $z \in \mathcal{V}$, if $z \in \mathcal{S}$ then
    $\sigma^{\mathcal{S}}_{pat^+}(z) = z_{pat^+} \notin \mathcal{V}$, else
    $\sigma^{\mathcal{S}}_{pat^+}(z) = z \in \mathcal{V}$.
    Therefore, $x \in \mathcal{S}$ if and only if $y \in \mathcal{S}$.
    We proceed by case analysis:
    \begin{itemize}
        \item Suppose $x \in \mathcal{S}$ and $y \in \mathcal{S}$.
              Therefore, $\sigma^{\mathcal{S}}_{pat^+}(x) = x_{pat^+}$ and
              $\sigma^{\mathcal{S}}_{pat^+}(y) = y_{pat^+}$. The hypothesis $x_{pat^+} = y_{pat^+}$
              implies that $x = y$.

        \item Suppose $x \notin \mathcal{S}$ and $y \notin \mathcal{S}$.
              Therefore, $\sigma^{\mathcal{S}}_{pat^+}(x) = x$ and
              $\sigma^{\mathcal{S}}_{pat^+}(y) = y$. The equality hypothesis leads to $x = y$.
    \end{itemize}
    Then we have $\sigma^{\mathcal{S}}_{pat^+}(x) = \sigma^{\mathcal{S}}_{pat^+}(y)$ that implies
    $x = y$. The other direction of the equivalence is verified, we can say that
    $\sigma^{\mathcal{S}}_{pat^+}$ is a renaming function on $\mathcal{V}$.

    Let $pat^+_1, pat^+_2 \in \mathcal{P}$ and $x \in \mathcal{S}$ such that
    $\sigma^{\mathcal{S}}_{pat_1^+}(x) = \sigma^{\mathcal{S}}_{pat_2^+}(x)$.
    %
    We know that $\sigma^{\mathcal{S}}_{pat_{1,2}^+}(x) = x_{pat_{1,2}^+}$, therefore the equality
    hypothesis implies that $x_{pat_{1}^+} = x_{pat_{2}^+}$ which implies that $pat_1^+$ equals
    $pat_2^+$.
    %
    Then the set $\Sigma^{\mathcal{S}}_{\mathcal{P}}$ is disjoint on $\mathcal{S}$.
\end{skippableproof}

\paragraph{Dependency relation}
\label{sec:semantics:grsync:causality:dependencies:rel}

\newcommand{\depRel}[3]{{\ensuremath{#2 \, {\gtrdot}_{\mathtiny{#1}} \, #3}}}
\newcommand{\depOrd}[3]{{\ensuremath{#2 \, {>}_{\mathtiny{#1}} \, #3}}}

We want to define a relation between identifier that represents the fact that an identifier has to
be evaluated before an other identifier, because the latter depends on the first one.

\begin{definition}[Dependency relation]
    Let $eq^+$ be a set of \GRSync equations, and $x$ and $y$ two identifiers.
    The relation $\depRel{eq^+}{x}{y}$, defined in
    \Crefrange{rule:semantics:grsync:causality:dependencies:rel:eq}%
    {rule:semantics:grsync:causality:dependencies:rel:eq:when}, holds when $x$ depends on $y$ in $eq^+$,
    meaning that the computation of $x$ needs the computation of $y$ to be performed.
\end{definition}

\subparagraph{Blocks of equations}

The dependency relation over a set of equations $eq^+$ holds whenever it holds for one of its
individual equations. If there exists an equation $eq$ in $eq^+$ where the dependency relation holds
between identifiers $x$ and $y$, then it holds for the entire set $eq^+$.

\begin{equation}
    \depRel{eq^+}{}{} = \bigcup \depRel{eq}{}{}^+ \label{rule:semantics:grsync:causality:dependencies:rel:eq}
\end{equation}

\subparagraph{Declarations}

In the case of a declaration of the form $x^+ = e$, the identifier $x$ depends on the identifiers
required to evaluate the expression $e$.

\begin{equation}
    \depRel{x^+ = e}{}{} = \{x^+\} \times \deps{e} \label{rule:semantics:grsync:causality:dependencies:rel:eq:decl}
\end{equation}

\subparagraph{Match}

In a \kwf{match} equation of the form:
$$\match{e}{x^+}{(\branch{pat_i^+}{e_i}{eq_i^+})^+},$$
the dependency relation is determined by various elements.
%
The identifiers $x^+$ being defined by the \kwf{match} equation depend on the expression $e$ that
is evaluated to determine the matching branch and on the guard $e_i$ of every branch.
%
For each branch of pattern $pat_i^+$, the equation block $eq_i^+$ introduces local dependencies that
have to be considered. However, as previously discussed in Paragraph~\ref{sec:semantics:grsync:causality:dependencies:renaming},
a simple union of those dependencies will lead to false causality breaches. The local dependencies
have to be renamed using the set of renaming functions
$\{\sigma^{x^+}_{pat^+} \mid pat^+ \in \{(pat_i^+)^+\}\}$, defined in
\Cref{def:semantics:grsync:causality:dependencies:renaming}, disjoint on $x^+$
(\Cref{lem:semantics:grsync:causality:dependencies:renaming}), that maps an identifier $x$ defined by the
\kwf{match} equation to a labelled version $x_{pat^+}$ and other identifiers to themselves.
%
The globally defined identifiers also depend on the branch-local ones.

\begin{equation}
    \begin{split}
        \depRel{\matchTab{e}{x^+}{(\branch{pat_i^+}{e_i}{eq_i^+})^+}}{}{} =
         & \{x^+\} \times \left( \deps{e}
        \cup \bigcup \deps{e_i}^+
        \cup \bigcup \sigma^{x^+}_{pat_i^+}\left(\{x^+\}\right)\right)         \\
         & \cup \bigcup \sigma^{x^+}_{pat_i^+}\left(\depRel{eq_i^+}{}{}\right)
    \end{split}
    \label{rule:semantics:grsync:causality:dependencies:rel:eq:match}
\end{equation}

\subparagraph{When}

The dependency relation for \kwf{when} constructs follows a similar pattern.
%
The event patterns $epat_i^+$ used in the branches introduce dependencies on the
identifiers involved in those patterns, as they determine when the branch is activated. In the
presence of guard expressions $e_i$, identifiers $x^+$ defined by the \kwf{when} construct also
depend on these guards, as the guards must be evaluated before proceeding with the corresponding
branch.
%
Similarly to the \kwf{match} equation, the blocks of equations $eq_i^+$ within the branches
contribute to the dependency relation.
%
The dependency relation also accounts for memory of signals, that can be used in all branches. For
all signal identifier $y$ defined by the \kwf{when} equation and all event pattern used in its
branches $epat^+$, $y$ and $\sigma^{x^+}_{epat^+}(y)$ depend on $\last{y}$.
%
Similarly to the \kwf{match} equation, the globally defined identifiers depend on the branch-local ones.

\begin{equation}
    \begin{split}
        \depRel{\whenTab{x^+}{x_m^+ = c^+}{(\branch{epat_i^+}{e_i}{eq_i^+})^+}}{}{} =
         & \{x^+\} \times \left( \bigcup \deps{epat_i}^+
        \cup \bigcup \deps{e_i}^+
        \cup \bigcup \sigma^{x^+}_{epat_i^+}\left(\{x^+\}\right)\right)         \\
         & \cup \left(\{x_m^+\}
        \cup \bigcup \sigma^{x_m^+}_{epat_i^+}\left(\{x_m^+\}\right)\right)
        \times \left\{(\last{x_m})^+\right\}                                    \\
         & \cup \bigcup \sigma^{x^+}_{epat_i^+}\left(\depRel{eq_i^+}{}{}\right)
    \end{split}
    \label{rule:semantics:grsync:causality:dependencies:rel:eq:when}
\end{equation}

\paragraph{Dependency order}
\label{sec:semantics:grsync:causality:dependencies:depOrd}

The relation \depRel{eq^+}{}{} only represents direct dependencies in $eq^+$. We can extend this
relation with transitivity, which defines the relation \depRel{eq^+}{}{} representing an order in
which identifiers of $eq^+$ should be computed.

\begin{definition}[Dependency order]
    Let $eq^+$ be a set of \GRSync equations.
    The transitive closure of the relation \depRel{eq^+}{}{} is denoted by \depOrd{eq^+}{}{}.
\end{definition}

\subsubsection{Causality Property}
\label{sec:semantics:grsync:causality:causal}

The causality of a component is a key predicate in \Crefrange{thm:semantics:grsync:theorems:productivity}{thm:semantics:grsync:theorems:equivalence},
as it states that no identifier from the component depends on itself, so that a mechanized
computation of the outputs from the inputs is possible.
%
The causality property is formally defined in \Cref{def:semantics:grsync:causality:causal:prop}.

\begin{definition}[Causality property]
    \label{def:semantics:grsync:causality:causal:prop}
    Let $eq^+$ be a set of equations. The equations $eq^+$ are \emph{causal} \iff there is no
    identifier $x$ such that \depOrd{eq^+}{x}{x}, or in other words \iff \depOrd{eq^+}{}{} is a
    strict partial order.

    A \component{f}{x^+_i}{x^+_o}{eq^+}{x^+_m = c^+} is \emph{causal} \iff its equations $eq^+$ are
    causal.
\end{definition}

Later in this section, we will prove that the causality of a component implies the existence of
scheduling (\Cref{prop:semantics:grsync:causality:scheduling}) in which equations can be computed
without any deadlock caused by circular dependencies.

\subsubsection{Rank}
\label{sec:semantics:grsync:causality:rank}

\newcommand{\rank}[2]{{\ensuremath{\textit{rk}_{#1}(#2)}}}

In a causal component, identifiers can be classified by dependency depth, that is the maximum length
of their dependency sequence. We call it the rank of the identifier and define it as follows.

\begin{definition}[Rank]
    \label{def:semantics:grsync:causality:rank}
    Let $eq^+$ be a causal set of equations. If $x$ is an identifier, $\rank{eq^+}{x}$ denotes the
    rank of $x$ and is defined as follows:
    \begin{equation*}
        \rank{eq^+}{x} = \text{max}\left(\left\{n \left| \begin{array}{l}
                x_0, \dots, x_n \in \defs{eq^+} \cup \deps{eq^+}, \\
                \qquad x_0 = x \, \wedge \, \forall k \in [0, n - 1], \, \depOrd{eq^+}{x_k}{x_{k+1}}
            \end{array}\right.\right\}\right)
    \end{equation*}
\end{definition}

\begin{lemma}[Rank well-defined]
    \label{lem:semantics:grsync:causality:rank:well_def}
    Let $eq^+$ be a causal set of equations and $x$ be an identifier.
    The rank $\rank{eq^+}{x}$ is well-defined.
\end{lemma}
\begin{skippableproof}
    Let $eq^+$ be a causal set of equations and $x$ be an identifier.
    %
    There exists at least one a simple sequence satisfying the elaboration of $\rank{eq^+}{x}$:
    $(x_0)$ with $x_0 = x$.
    %
    For all $n \in \N$ and sequences $x_0, \dots, x_n \in \defs{eq^+} \cup \deps{eq^+}$
    such that $x_0 = x$ and $\depOrd{eq^+}{x_k}{x_{k+1}}$ for all $k \in [0, n - 1]$, we have $n$
    bounded by $\left|\defs{eq^+} \cup \deps{eq^+}\right|$ (otherwise an identifier would appear
    twice in the sequence, creating a causality violation).
    %
    Therefore, the maximum size of such sequences exists and the rank $\rank{eq^+}{x}$ is
    well-defined.
\end{skippableproof}

\begin{lemma}[Rank decreases]
    \label{lem:semantics:grsync:causality:rank:decrease}
    Let $eq^+$ be a causal set of equations.
    For all identifiers $x$ and $y$ in $\defs{eq^+} \cup \deps{eq^+}$:
    $$\depOrd{eq^+}{x}{y} \Rightarrow \rank{eq^+}{x} > \rank{eq^+}{y}.$$
\end{lemma}
\begin{skippableproof}
    Let $eq^+$ be a causal set of equations, and $x, y \in \defs{eq^+} \cup \deps{eq^+}$ such that
    $\depOrd{eq^+}{x}{y}$.
    %
    \Cref{lem:semantics:grsync:causality:rank:well_def} ensures there exists a sequence
    $(y_0, \dots, y_n) \in \defs{eq^+} \cup \deps{eq^+}$ such that $\rank{eq^+}{y} = n$, $y_0 = y$,
    and $\depOrd{eq^+}{y_k}{y_{k+1}}$ holds for all $k \in [0, n - 1]$. We have
    $$
        (x, y_0, \dots, y_n) \in \left\{(z_0, \dots, z_m) \mid
        z_0 = x \,
        \wedge \, \forall k \in [0, m - 1], \, \depOrd{eq^+}{z_k}{z_{k+1}}
        \right\},
    $$
    which implies that $\rank{eq^+}{x} \ge n+1$, and finally $\rank{eq^+}{x} > \rank{eq^+}{y}$.
\end{skippableproof}

\begin{lemma}[Bounded]
    \label{lem:semantics:grsync:causality:rank:bounded}
    Let $eq^+$ be a causal set of equations.
    For all identifier $x \in \defs{eq^+} \cup \deps{eq^+}$, we have $\rank{eq^+}{x} \le |eq^+|$.
\end{lemma}
\begin{skippableproof}
    Let $eq^+$ be a causal set of equations. Let $x \in \defs{eq^+} \cup \deps{eq^+}$ be an
    identifier such that $r = \rank{eq^+}{x}$.
    %
    Therefore there exist a sequence $(x_0, \dots, x_r) \in \defs{eq^+} \cup \deps{eq^+}$ such that
    $x_0 = x$ and $\forall k \in [0, n - 1], \, \depOrd{eq^+}{x_k}{x_{k+1}}$.

    All identifiers in the sequence $(x_0, \dots, x_r)$ are different (otherwise $eq^+$ would not be
    causal), and all identifiers in $(x_0, \dots, x_{r-1})$ are defined in $eq^+$ (otherwise there
    would be no dependency relation between them).\footnote{The identifier $x_r$ might not be
        defined by $eq^+$, we just know that $x_{r-1}$ depends on it in $eq^+$ which implies that
        $x_{r-1}$ is defined in those equations.}

    The dependency relation induced by an equation is identical for all the identifiers it defines.
    Therefore, there cannot be two identifiers of this sequence that are defined by the same
    equation $eq' \in eq^+$, as they would introduce a causality loop.

    Therefor, all identifiers in the sequence $(x_0, \dots, x_r)$ are from different equations.
    It implies that $r \le |eq^+|$.
\end{skippableproof}

\subsubsection{Induction}
\label{sec:semantics:grsync:causality:induction}

Reasoning about \GRSync components often requires establishing properties for all identifiers
defined or used in a component. Since the dependency relation in causal components is acyclic, we
can process identifiers in an order that respects their dependencies. This approach motivates an
induction principle based on the structure of these dependencies.

This idea is formalized in the \emph{causal induction}~\cite{bourke:hal-03936656}, where a property
is proven for an identifier by assuming it holds for all its dependencies. In this manuscript, we
use an equivalent induction principle based on the \emph{rank} of identifiers, as previously
defined.

\begin{definition}[Rank induction]
    \label{def:semantics:grsync:causality:induction:rank}
    Let $eq^+$ be a causal set of equations, and let $S = \defs{eq^+} \cup \deps{eq^+}$ denote the
    set of all relevant identifiers for the set of equations.
    Let \P{x} be a property on the identifier $x \in S$.
    Suppose:
    \begin{itemize}
        \item \textit{Base case:} \P{x} holds for all $x \in S$ with $\rank{eq^+}{x} = 0$.
        \item \textit{Inductive step:} For any $n > 0$, if \P{y} holds for all $y \in S$ with
              $\rank{eq^+}{y} < n$, then \P{x} holds for all $x \in S$ with $\rank{eq^+}{x} = n$.
    \end{itemize}
    Then \P{x} holds for all $x \in S$.
\end{definition}

\begin{lemma}[Soundness of rank induction]
    \label{lem:semantics:grsync:causality:induction:soundness}
    The rank induction principle is sound.
\end{lemma}
\begin{skippableproof}
    By \Cref{lem:semantics:grsync:causality:rank:well_def} and
    \Cref{lem:semantics:grsync:causality:rank:bounded}, for any causal set of equations $eq^+$, the
    rank function \rank{eq^+}{x} assigns a natural number to each identifier $x \in S$, and the set
    of possible ranks is finite.
    %
    The set of natural numbers \N is well-founded under the usual order $<$, meaning it has no
    infinite descending chains. The induction principle stated above is thus a direct application
    of well-founded induction on \N.

    Since the base case holds for rank $0$, and the inductive step ensures the property for rank
    $n$ based on all smaller ranks, it follows by induction that \P{x} holds for all $x \in S$.
\end{skippableproof}

\subsubsection{Scheduling}
\label{sec:semantics:grsync:causality:scheduling}

Causality in a set of equations guarantees that there are no cyclic dependencies among identifiers.
\Cref{prop:semantics:grsync:causality:scheduling} states that a valid scheduling exists: there is an
order in which the identifiers can be defined such that, for each one, all of its dependencies are
resolved beforehand.

\begin{proposition}
    \label{prop:semantics:grsync:causality:scheduling}
    Let $eq^+$ be a causal set of equations. There exists a
    scheduling $(x_0, \dots, x_n)$ such that $\{x_0, \dots, x_n\} = \defs{eq^+} \cup \deps{eq^+}$,
    and for all $k \in [0, n]$ and $y \in \defs{eq^+} \cup \deps{eq^+}$ such that
    $\depOrd{eq^+}{x_k}{y}$, we have $y \in \{x_0, \dots, x_{k-1}\}$.
\end{proposition}
\begin{skippableproof}
    Let $eq^+$ be a causal set of equations. We will schedule the identifiers based on their rank in
    increasing order.
    %
    \Cref{lem:semantics:grsync:causality:rank:well_def} ensures that all
    $x \in \defs{eq^+} \cup \deps{eq^+}$ has a rank and then will be scheduled.
    %
    \Cref{lem:semantics:grsync:causality:rank:bounded} states that the rank is bounded, so that we
    can define $R = \text{max}(\{\rank{eq^+}{x} \mid x \in \defs{eq^+}\})$.
    %
    We can then classify the identifiers defined in $eq^+$ by their rank and establish the sets
    $X_r = \{x \mid x \in \defs{eq^+} \wedge \rank{eq^+}{x} = r\}$ for all $r \in [0, R]$.
    %
    \Cref{lem:semantics:grsync:causality:rank:decrease} ensures that for all rank $r \in [0, R]$,
    identifier $x \in X_r$, and identifier $y \in \defs{eq^+} \cup \deps{eq^+}$ such that
    $\depOrd{eq^+}{x}{y}$, we have $\rank{eq^+}{x} > \rank{eq^+}{y}$. It implies that $y$ is
    scheduled before $x$ in the scheduling $(X_{r \in [0, R]})$.
\end{skippableproof}

\begin{remark}
    This scheduling methodology is an implementation of the topological sorting on the graph whose
    nodes are the identifiers and the edges are the dependency relation \depRel{eq^+}{}{}.
    %
    Another implementation, using another notion of rank that loops on the direct computation of the
    scheduling order, can be found in \cite{izerrouken:hal-02270296}.
\end{remark}
